\begin{abstract}
Password-authenticated searchable encryption (PASE) lets a user store encrypted files on cloud servers and search
them by keyword while remembering only a password. QPASE, by Jiang and Wang, rebuilds PASE on lattices so that it
resists quantum attacks~\cite{jiang2024qpase}. We implemented QPASE in Python from the paper alone. It passes
\UnitTests{} unit tests, recovers the same user key in \LoginsIdentical{} logins and returns exactly the matching
documents. Building it showed that the lattice PRF sets the cost of every user-side phase, that the paper's 128-bit
parameter set allows offline password guessing in the ring $\mathbb{Z}_q[X]/(X^{595}+1)$, and that the linear
key-update commitments reveal the shares. We then propose a counter-based search index that keeps the search key
on the user's device and gives forward privacy. At the slowest server it searched \ScalingSpeedupMin{} to
\ScalingSpeedupMax{} times faster than the paper's scheme for $10^2$ to $10^5$ stored files, mostly because the
server no longer evaluates the PRF for each match. The user's cost is unchanged.

\medskip
\noindent\textbf{Keywords:} searchable encryption, password authentication, lattice cryptography, forward privacy,
post-quantum security.
\end{abstract}
